\documentclass[aspectratio=169, 12pt]{beamer}

% --- Theme ---
\usetheme{metropolis}
\usecolortheme{default}
\usepackage[utf8]{inputenc}
\usepackage[T1]{fontenc}
\usepackage{booktabs}
\usepackage{graphicx}
\usepackage{hyperref}

% --- Colours ---
\definecolor{synblue}{HTML}{1f6aa5}
\definecolor{syngreen}{HTML}{2a9d5c}
\definecolor{synorange}{HTML}{e76f1a}
\setbeamercolor{frametitle}{bg=synblue}
\setbeamercolor{progress bar}{fg=syngreen}

% --- Title ---
\title{Patents, Papers, Parts \& Planet}
\subtitle{Tracing Synthetic Biology Innovation Trajectories\\from Student Projects to Patents}
\author{[Author Name]}
\institute{CSH Vienna / TU Straubing}
\date{\today}

% ============================================================
\begin{document}
% ============================================================

\maketitle

% ------------------------------------------------------------
\begin{frame}{Overview}
  \tableofcontents
\end{frame}

% ============================================================
\section{Motivation}
% ============================================================

\begin{frame}{Why study innovation trajectories in synthetic biology?}
  \begin{itemize}
    \item Synthetic biology is a frontier technology with real-world impact
    \item Innovation leaves traces across \textbf{multiple knowledge artifacts}:
      student projects, papers, patents
    \item These traces cluster \textbf{geographically}
    \item We want to understand: do these artifact types form related
      \emph{local innovation trajectories}?
  \end{itemize}
\end{frame}

% ============================================================
\section{Research Question}
% ============================================================

\begin{frame}{Research question}
  \begin{block}{}
    \Large Do iGEM student projects, academic papers, and patents in synthetic
    biology occupy \textbf{semantically related regions} of the same knowledge
    space at the \textbf{city level}?
  \end{block}

  \vspace{1em}
  \textbf{Working hypothesis:} a plausible innovation sequence
  \begin{enumerate}
    \item Student projects introduce early ideas
    \item Academic papers develop them
    \item Patents reflect downstream application
  \end{enumerate}

  \vspace{1em}
  \textit{We test association, not causation.}
\end{frame}

% ============================================================
\section{Data}
% ============================================================

\begin{frame}{Three types of artifact}
  \begin{columns}
    \column{0.33\textwidth}
    \textbf{\textcolor{synblue}{Papers}}\\
    OpenAlex API\\
    \textasciitilde{}[N] records\\
    [year range]

    \column{0.33\textwidth}
    \textbf{\textcolor{synorange}{Patents}}\\
    Lens.org API\\
    \textasciitilde{}[N] records\\
    [year range]

    \column{0.33\textwidth}
    \textbf{\textcolor{syngreen}{iGEM Projects}}\\
    iGEM Registry\\
    \textasciitilde{}[N] teams\\
    [year range]
  \end{columns}

  \vspace{2em}
  All normalized to a \textbf{shared schema}: id, type, title, text, year, city, country, \ldots
\end{frame}

% ============================================================
\section{Methods}
% ============================================================

\begin{frame}{Pipeline overview}
  \begin{enumerate}
    \item \textbf{Ingest} papers, patents, iGEM projects from APIs
    \item \textbf{Embed} each artifact's title + abstract with a sentence-transformer
    \item \textbf{Reduce} to 2D with UMAP
    \item \textbf{Cluster} with HDBSCAN
    \item \textbf{Geocode} institution locations
    \item \textbf{Tag} carbon-capture artifacts with keyword matching
  \end{enumerate}
\end{frame}

\begin{frame}{Embedding model}
  \begin{itemize}
    \item \texttt{all-MiniLM-L6-v2} (sentence-transformers library)
    \item 384-dimensional dense vectors
    \item Captures semantic meaning of title + abstract
    \item Fast, local, no API key required
  \end{itemize}

  \vspace{1em}
  \textit{Reimers \& Gurevych (2019), Sentence-BERT, EMNLP 2019}
\end{frame}

\begin{frame}{Dimensionality reduction and clustering}
  \begin{columns}
    \column{0.5\textwidth}
    \textbf{UMAP} (2D projection)
    \begin{itemize}
      \item Preserves local and global structure
      \item cosine metric
      \item $n\_\text{neighbors} = 15$
    \end{itemize}

    \column{0.5\textwidth}
    \textbf{HDBSCAN} (clustering)
    \begin{itemize}
      \item No fixed number of clusters
      \item Handles noise/outliers
      \item $\text{min\_cluster\_size} = 10$
    \end{itemize}
  \end{columns}
\end{frame}

% ============================================================
\section{Results}
% ============================================================

\begin{frame}{Semantic space — all artifact types}
  \begin{center}
    \textit{[Figure: UMAP scatter plot colored by artifact type]}\\
    \vspace{1em}
    \small Papers, patents, and iGEM projects in the same 2D space.
  \end{center}
\end{frame}

\begin{frame}{Cluster composition}
  \begin{center}
    \textit{[Figure: bar charts showing artifact type mix per cluster]}
  \end{center}
\end{frame}

\begin{frame}{City-level patterns}
  \begin{center}
    \textit{[Figure: map of synthetic biology activity by city]}
  \end{center}
\end{frame}

% ============================================================
\section{Case Study: Carbon Capture}
% ============================================================

\begin{frame}{Carbon capture in synthetic biology}
  \begin{itemize}
    \item Engineering microorganisms to fix CO\textsubscript{2} or produce carbon-neutral fuels
    \item Tagged by keyword matching: [N] papers, [N] patents, [N] iGEM projects
    \item \textbf{Do they cluster together?}
  \end{itemize}
\end{frame}

\begin{frame}{Carbon capture — semantic view}
  \begin{center}
    \textit{[Figure: UMAP with carbon-capture artifacts highlighted]}
  \end{center}
\end{frame}

\begin{frame}{Carbon capture — timeline}
  \begin{center}
    \textit{[Figure: year-by-year counts of papers, patents, projects]}
  \end{center}
\end{frame}

% ============================================================
\section{Discussion}
% ============================================================

\begin{frame}{What we found}
  \begin{itemize}
    \item [Main finding 1]
    \item [Main finding 2]
    \item [Main finding 3]
  \end{itemize}
\end{frame}

\begin{frame}{Limitations}
  \begin{itemize}
    \item Keyword-based corpus construction can miss or over-include
    \item Geocoding is approximate; some locations missing
    \item Semantic similarity $\neq$ citation or collaboration link
    \item No causal claim — observational association only
  \end{itemize}
\end{frame}

% ============================================================
\section{Implications}
% ============================================================

\begin{frame}{Implications}
  \begin{itemize}
    \item iGEM as an early-signal dataset for innovation monitoring
    \item Synthetic biology subfields may have detectable local trajectories
    \item Shared semantic space as a tool for cross-artifact analysis
    \item Carbon capture: [substantive finding about this subfield]
  \end{itemize}
\end{frame}

\begin{frame}[standout]
  Thank you.\\
  \vspace{1em}
  \small
  Code and data: \href{https://github.com/Zer0Juice/synbio-diversification}{github.com/Zer0Juice/synbio-diversification}\\
  Website: [GitHub Pages URL]
\end{frame}

% ============================================================
\end{document}
